\section{The exact geometric multiplicity monoid over a prime field}

This section completes the geometric rational essential-image test
from the prime field in every slope sector.  It is a consequence of
the inspected Milne classification and the transition formula
CC-AG09.  The category is still the rational numerical category
generated by abelian varieties, finite étale schemes and Tate
twists.  No statement about the missing cross-correspondences of
an arbitrary variety is used or implied.

\subsection{CC-AG13: a computable generator in all slopes}

Fix a prime \(p\).  Let \(\xi\in\mathcal W_p\) have weight
\(w\in\mathbb Z\), and choose a stabilized representative
\[
 \xi=[\pi,b],\qquad b\ge1,\qquad
 K=K_\xi=\mathbb Q(\pi^N)\quad\text{for every }N\ge1.
                                                        \tag{CC-AG.80}
\]
Write \(d_\xi\) for the geometric Schur index of its rational
simple object \(S_\xi\).  In particular, its realization has
dimension \(d_\xi[K:\mathbb Q]\).
Define the exponent subgroup directly by
\[
 \mathscr E_\xi=
 \{n\in\mathbb Z:\ \text{there exists }\gamma\in K^\times
                    \text{ with }\gamma^b/\pi^n\in\mu(K)\}.
                                                        \tag{CC-AG.81}
\]
This definition uses no integrality assumption on the normalized
\(p\)-slopes and no fractional ideal with fractional exponents.
Then \(\mathscr E_\xi=h_\xi\mathbb Z\), where \(h_\xi>0\) and
\(h_\xi\mid b\), and
\[
 h_\xi=\min\{n>0:n\mid b,\ 
            X^b-\zeta\pi^n\text{ has a root in }K
                          \text{ for some }\zeta\in\mu(K)\}.
                                                        \tag{CC-AG.82}
\]
The positive integer \(h_\xi\) is independent of the chosen
stabilized representative: equivalently, it is the least positive
exponent \(n\) at which \(\xi\) has a representative in \(K\).

The geometric index divides this exponent:
\[
                              d_\xi\mid h_\xi.
                                                        \tag{CC-AG.83}
\]
The exact multiplicity monoid for base extension from
\(\mathbb F_p\) to \(\overline{\mathbb F}_p\) is
\[
                 H_{\mathbb F_p}(O_\xi)
                      =\frac{h_\xi}{d_\xi}\,\mathbb N.
                                                        \tag{CC-AG.84}
\]
Thus a rational geometric object
\[
                         \bigoplus_\xi S_\xi^{\oplus r_\xi}
\]
with finitely many distinct Galois orbits of labels descends
to \(\mathcal A_p\) if and only if
\[
                        h_\xi/d_\xi\mid r_\xi
                            \quad\text{for every }\xi.
                                                        \tag{CC-AG.85}
\]
Both integers in this test are effectively determined by the
finite equations (CC-AG.82) and the local-invariant formula
of CC-AG08.

If \(K\) has no real embeddings, choose a witness \(\gamma\)
for \(h_\xi\) in (CC-AG.81).  Every root
\[
                             \alpha^{h_\xi}=\gamma
                                                        \tag{CC-AG.86}
\]
is a weight-\(w\) Weil \(p\)-number with germ \(\xi\),
has centre degree
\[
                           [\mathbb Q(\alpha):K]=h_\xi,
                                                        \tag{CC-AG.87}
\]
and has source Schur index one.  Its rational simple object
base changes to \(S_\xi^{\oplus h_\xi/d_\xi}\), realizing the
least positive multiplicity by a single simple object.

\paragraph{Complete proof.}
Multiplication and inversion of witnesses in (CC-AG.81) show
that \(\mathscr E_\xi\) is a subgroup of \(\mathbb Z\).
It contains \(b\), with witness \(\pi\), so it is a nonzero
subgroup \(h_\xi\mathbb Z\), with \(h_\xi>0\) and \(h_\xi\mid b\).
Testing its positive divisors of \(b\) gives (CC-AG.82).

For \(n>0\), a witness satisfies
\(\gamma^b=\zeta\pi^n\).  At every complex embedding it
therefore has absolute value \(p^{nw/2}\); at every finite
place away from \(p\) its valuation is zero.  It is a
weight-\(w\) Weil \(p^n\)-number.  After raising the equation
to the order of \(\zeta\), it gives
\([\gamma,n]=[\pi,b]\).  Conversely, this equality of germs
means that \(\gamma^{bN}=\pi^{nN}\) for some \(N\ge1\).
Consequently \(\gamma^b/\pi^n\) is a root of unity in \(K\).
This proves the intrinsic representative description, and
hence independence in (CC-AG.82).

This is a finite exact test over a fixed number field.
For completeness, one explicit way to determine \(\mu(K)\)
is the bounded-coefficient construction already proved
in CC-AG11--12.  Put \(D=[K:\mathbb Q]\) and
\[
 R_D=\sum_{e=1}^{D}e\prod_{j=1}^{e}
                  \left(2\binom ej+1\right),\qquad
 N_D=\operatorname{lcm}(1,\ldots,R_D).
\]
The minimal polynomial of every power of a root of unity
in \(K\) has degree at most \(D\) and its coefficient of
degree \(e-j\) has absolute value at most \(\binom ej\).
There are at most \(R_D\) possible roots of all these
bounded polynomials.  Hence the order of every root of
unity in \(K\) is at most \(R_D\), and \(\mu(K)\) is precisely
the set of roots in \(K\) of \(X^{N_D}-1\).  Factoring this
polynomial and the finitely many polynomials in
(CC-AG.82) terminates.  This deliberately coarse bound is
used only for finiteness, not for computational efficiency.

First suppose that \(K\) has no real embeddings.
Let \(\gamma\) witness exponent \(h=h_\xi\), and let
\(v\mid p\) be a place of \(K\), with normalized integral
valuation \(\operatorname{ord}_v\), ramification index \(e_v\)
and residue degree \(f_v\).  The geometric local invariant,
computed using the representative \([\gamma,h]\), is
\[
 -\frac{\operatorname{ord}_v(\gamma)}
              {h e_v}[K_v:\mathbb Q_p]
       =-\frac{\operatorname{ord}_v(\gamma)f_v}{h}
                                             \quad\bmod\mathbb Z.
                                                        \tag{CC-AG.88}
\]
Its denominator divides \(h\).  There are no real-place
invariants, and all remaining invariants are zero.
The Schur index is their least common denominator, so
\(d_\xi\mid h\), proving (CC-AG.83) in this case.
This uses the geometric-Frobenius convention in the
corrected Milne manuscript, p.~17 and Theorems~2.15--2.16
and 2.18--2.19, pp.~18 and 20.

Take any source Weil \(p\)-number \(\alpha\) whose germ
orbit is \(O_\xi\).  After conjugation, assume
\([\alpha,1]=[\pi,b]\), and put \(L=\mathbb Q(\alpha)\).
For some \(N\ge1\),
\[
                              \alpha^{bN}=\pi^N.
                                                        \tag{CC-AG.89}
\]
The stabilization in (CC-AG.80) gives \(K\subset L\).
The element \(\theta=\alpha^b/\pi\) lies in \(\mu(L)\).
Set \(d=[L:K]\).  Taking the relative field norm gives
\[
 \frac{N_{L/K}(\alpha)^b}{\pi^d}
                         =N_{L/K}(\theta)\in\mu(K).
                                                        \tag{CC-AG.90}
\]
The norm of a root of unity is a root of unity because
the norm is multiplicative.  No Galois hypothesis on
\(L/K\) is required.  Thus \(d\in\mathscr E_\xi\), and
\[
                                     h\mid d.
                                                        \tag{CC-AG.91}
\]

The field \(L\) has no real embedding, since one would
restrict to a real embedding of \(K\).  At a place
\(z\mid p\) of \(L\), the local invariant of the source
simple over the prime field is
\[
 -\frac{\operatorname{ord}_z(\alpha)}
                 {\operatorname{ord}_z(p)}[L_z:\mathbb Q_p]
               =-\operatorname{ord}_z(\alpha)f_z
                                             \quad\bmod\mathbb Z.
                                                        \tag{CC-AG.92}
\]
It is zero.  The other invariants are also zero, and
there are no real places; hence the source index
\(d_\alpha\) is one.  By CC-AG09 its base-change
multiplicity is
\[
                   [L:K]\frac{d_\alpha}{d_\xi}
                         =\frac{d}{d_\xi}
                    \in\frac{h}{d_\xi}\mathbb N.
                                                        \tag{CC-AG.93}
\]
This proves the necessary divisibility for every
source simple and thus for every positive direct sum.

For attainment, take a root as in (CC-AG.86).
All complex absolute values are \(p^{w/2}\), and all
valuations away from \(p\) vanish, so \(\alpha\in W(p)\).
The exact equation gives
\([\alpha,1]=[\gamma,h]=\xi\).
Every positive power of \(\gamma\) generates \(K\).
Indeed, if \(\gamma^b=\zeta\pi^h\), raising to the
order of \(\zeta\) and then to any positive power
shows that the field of that power of \(\gamma\)
contains a positive power of \(\pi\), which generates
\(K\); the reverse inclusion is immediate.
In particular \(K\subset\mathbb Q(\alpha)\).
The degree of this extension is at most \(h\), by
(CC-AG.86), and is a positive multiple of \(h\), by
(CC-AG.91).  It is exactly \(h\).
Equation (CC-AG.92) gives source index one, so
(CC-AG.93) realizes \(h/d_\xi\).
Direct sums realize all its nonnegative multiples.
This proves (CC-AG.84) when \(K\) has no real embeddings.

It remains to handle every possible real-place case.
If \(K\) has a real embedding \(\sigma\), then
\(\sigma(\pi)^2=p^{bw}\) by the Weil absolute-value
condition.  Injectivity of \(\sigma\) gives
\(\pi^2=p^{bw}\).  Stabilization therefore implies
\[
                             K=\mathbb Q(\pi^2)=\mathbb Q.
                                                        \tag{CC-AG.94}
\]
The germ is uniquely determined by \(w\); signs are
roots of unity and disappear in the germ.
If \(w\) is even, it has the stage-one representative
\(p^{w/2}\in\mathbb Q\), so \(h=1\).  Its real and
\(p\)-adic invariants are integral, giving \(d_\xi=1\).
The corresponding rational one-dimensional Tate
source realizes multiplicity one.

If \(w\) is odd, a rational representative at exponent
\(n\) would have absolute value \(p^{nw/2}\).
This is rational precisely when \(n\) is even.
The representative \(p^w\) at exponent two exists,
including when \(w<0\); hence \(h=2\).
The real invariant \(w/2\) and the \(p\)-adic invariant
\(-w/2\) have denominator two, so \(d_\xi=2\).
Choose instead the equivalent stage-two representative
\(\gamma=-p^w\), and choose
\[
                                  \alpha^2=-p^w.
                                                        \tag{CC-AG.95}
\]
Its field is imaginary quadratic, because \(p^w>0\)
and \(w\) is odd; its degree over \(\mathbb Q\) is two.
It is a weight-\(w\) Weil \(p\)-number with germ \(\xi\).
Equation (CC-AG.92) and absence of real places give
source index one.  The transition multiplicity is
\(2/2=1=h/d_\xi\).
Every nonnegative integer is now realized.  Other
source labels with real centre may have index two;
this does not affect the exact monoid \(\mathbb N\).
Thus (CC-AG.83)--(CC-AG.84) hold also in the real case.

Finally, simple rational objects are classified by
Galois orbits of labels.  Base extension of one simple
has only the one target orbit occurring in CC-AG09.
Semisimple decomposition and positivity therefore
separate distinct target orbits.  Applying the proved
single-orbit criterion to each proves (CC-AG.85).
The zero object has only zero multiplicities and is
included. \(\square\)

\paragraph{Scope relative to the preceding results.}
CC-AG11 gives \(h/\gcd(h,a)\) in the integral-slope sector
over every finite source field.  CC-AG13 gives the entire
all-slope geometric monoid over the prime field, with the
geometric index \(d_\xi\) retained.  Over
\(\mathbb F_{p^a}\) for \(a>1\), the source invariant is
\(-\operatorname{ord}_z(\alpha)f_z/a\); it need not vanish.
No all-slope cyclic formula for those source fields is
asserted here.  CC-AG12 still supplies the finite exact
decision and all preimage isomorphism classes for each
specified target object in that general case.

\subsection{Primary references and exact conventions}

\begin{itemize}
\item Milne, J.~S. (1994). Motives over finite fields. In
U.~Jannsen, S.~Kleiman, \& J.-P.~Serre (Eds.),
\emph{Motives} (Proceedings of Symposia in Pure Mathematics,
Vol.~55, Part~1, pp.~401--459). American Mathematical Society.
Corrected author version dated November~2, 2019:
\url{https://www.jmilne.org/math/xnotes/Seattle1.pdf}.
Geometric-Frobenius invariant convention, p.~17;
Theorems~2.15--2.16, p.~18; Theorems~2.18--2.19,
p.~20.  These inputs are applied only within the
established abelian--Artin--Tate category.
\item Milne, J.~S. (2020, July~19).
\emph{Algebraic number theory} (Version~3.08)
[Author course notes].
\url{https://www.jmilne.org/math/CourseNotes/ANT.pdf}.
Relative field norm, Proposition~2.19,
statement p.~31 and proof pp.~31--32;
bounded-polynomial root-of-unity argument,
Proposition~5.5 and Corollary~5.6, p.~87.
\end{itemize}
